% !PRINTED 346
\weberpage{346}
Wendet man die Formel (1) auf eine reelle Gleichung an, so ist $\alpha^m + \beta^m + \gamma^m + \cdots$ eine reelle Zahl; und wenn die absolut grösste Wurzel $\alpha$ reell ist, so ist auch die $m^{\mathrm{te}}$ Wurzel reell zu nehmen. Ist aber $m$ eine gerade Zahl, so muss man, um über das Vorzeichen zu entscheiden, noch wissen, ob $\alpha$ positiv oder negativ ist. N\"othigenfalls wird darüber durch Einsetzen des gefundenen Näherungswerthes in die gegebene Gleichung entschieden.

Ebenso wird man, wenn man die Formel (1) auf eine imaginäre Gleichung anwendet, entscheiden müssen, welche der verschiedenen $m^{\mathrm{ten}}$ Wurzeln die richtige Annäherung giebt.

Die Ausdrücke
\[
s_m = \alpha^m + \beta^m + \gamma^m + \cdots
\]
lassen sich besonders einfach dann berechnen, wenn man für $m$ die auf einander folgenden Potenzen von 2, also $m = 2, 4, 8, 16 \ldots$ setzt, und liefern dann sehr gute Resultate.

Es ist $s_2$ der negative Co\"efficient der $(n-1)^{\mathrm{ten}}$ Potenz der Unbekannten in der Gleichung, deren Wurzeln die Quadrate der Wurzeln von $f(x)$ sind. Diese Gleichung wird aber leicht auf folgende Weise gebildet.

Man fasse in $f(x)$ die Glieder mit geraden und mit ungeraden Potenzen von $x$ zusammen und setze demnach
\[
f(x) = \varphi(x^2) + x\,\psi(x^2).
\]
Es ist dann
\[
f_1(x) = \varphi(x)^2 - x\,\psi(x)^2 = 0
\]
die Gleichung, deren Wurzeln die Quadrate der Wurzeln von $f(x)$ sind. Behandelt man $f_1(x)$ ebenso, so erhält man eine Gleichung, deren Wurzeln die vierten Potenzen der Wurzeln von $f(x)$ sind u. s. f.

Die Ausführung dieser Rechnung ist sehr einfach und führt meist nach wenigen Schritten zu einer guten Näherung.

Wir betrachten einige Beispiele.

Zunächst nehmen wir die Gleichung
\weq{4}{x^3 - 2x - 2 = 0,}
die eine reelle positive und zwei imaginäre Wurzeln hat. Die reelle Wurzel ist grösser als 1,7, und da das Product aller drei Wurzeln gleich 2 und $(1,7)^3 > 2$ ist, so ist der absolute Werth der beiden imaginären Wurzeln kleiner als die reelle Wurzel. Also ist die reelle Wurzel die absolut grösste und die \gesperrt{Gräffe'}-
%CONT>>
